% =====================================================================
%  PCRN: A Pocket-derived Co-occurrence Residue Network for
%        Analyzing Protein Pocket Dynamics
%
%  Self-contained Overleaf project. Just upload this file (and the
%  references.bib in the same folder) to Overleaf and click "Recompile".
%  Engine: pdfLaTeX. Bibliography: BibTeX (run pdflatex -> bibtex ->
%  pdflatex -> pdflatex, or just press Recompile twice on Overleaf).
% =====================================================================

\documentclass[11pt,a4paper]{article}

% --------- Encoding & fonts ---------
\usepackage[utf8]{inputenc}
\usepackage[T1]{fontenc}
\usepackage{lmodern}

% --------- Layout ---------
\usepackage[margin=1in]{geometry}
\usepackage{setspace}
\onehalfspacing

% --------- Math & symbols ---------
\usepackage{amsmath,amssymb}

% --------- Graphics & color ---------
\usepackage{graphicx}
\usepackage{xcolor}

% --------- Hyperlinks ---------
\usepackage[colorlinks=true,
            linkcolor=blue!60!black,
            citecolor=blue!60!black,
            urlcolor=blue!60!black]{hyperref}

% --------- Section formatting ---------
\usepackage{titlesec}
\titleformat{\section}{\large\bfseries}{\thesection}{0.6em}{}
\titleformat{\subsection}{\normalsize\bfseries}{\thesubsection}{0.6em}{}

% --------- Better tables (if needed later) ---------
\usepackage{booktabs}

% --------- Bibliography style ---------
\usepackage[numbers,sort&compress]{natbib}
\bibliographystyle{unsrtnat}

% =====================================================================
%   Title block
% =====================================================================
\title{\textbf{PCRN: A Pocket-derived Co-occurrence Residue Network for
       Analyzing Protein Pocket Dynamics}}

\author{Anonymous Author(s) \\ \small Affiliation \\
        \small \texttt{email@institution.edu}}

\date{}

% =====================================================================
\begin{document}
% =====================================================================

\maketitle

% ---------------------------------------------------------------------
%   Abstract
% ---------------------------------------------------------------------
\begin{abstract}
\noindent
Protein pockets are inherently dynamic structures whose stability,
merging, splitting, and allosteric coupling underpin protein function
and drug binding. Existing residue-level networks predominantly encode
physical proximity or non-covalent interactions, capturing structural
contacts but overlooking the functional co-membership of residues
within transient pockets. We introduce \textbf{PCRN
(Pocket-derived Co-occurrence Residue Network)}, in which an edge
between two residues is weighted by how frequently they co-occur as
members of the same pocket across a conformational ensemble, with node
weights reflecting per-residue pocket participation frequency. The
construction is alignment-free, builds directly on per-frame pocket
detections from \textsc{fpocket}, and shifts the network's semantic
level from physical contact to functional co-membership. We show that
PCRN exposes pocket-relevant features that contact-based networks
miss: tightly co-occurring residue cliques mark stable pocket cores,
low-weight bridges identify hinges where pockets merge or split, and
community structure tracks the temporal organization of pocket
regions. PCRN is computationally lightweight, complementary to
existing residue interaction networks, and offers a generalizable
lens for studying dynamic pocket behavior.

\medskip
\noindent\textbf{Keywords:} protein pocket dynamics, residue
co-occurrence network, alignment-free analysis, molecular dynamics,
visual analytics, allosteric communication.
\end{abstract}

% =====================================================================
\section{Introduction}
% =====================================================================

\subsection{The Importance of Protein Pockets}
Protein pockets --- also referred to as cavities or binding sites ---
are geometrically concave regions formed by specific residues on the
protein surface or in its interior, and they play central roles in
ligand recognition, catalysis, and allosteric regulation
\citep{stank2016protein,krone2016visual}. In structural biology and
drug design, the accurate localization and characterization of pockets
serve as a critical starting point for hit discovery, virtual
screening, and mechanistic interpretation
\citep{meller2025cryptic,krivak2024comparative}. Increasing evidence,
however, indicates that many functionally relevant pockets ---
including cryptic pockets, allosteric sites, and transient pockets
--- are not visible in static crystal structures and can only be
revealed by examining the dynamic conformational behavior of the
protein \citep{stank2016protein,meller2025cryptic}. Consequently,
systematic analysis of pockets and their dynamics has become a central
topic in modern structure-based drug discovery.

\subsection{The Development of Static Pocket Detection Algorithms}
Pocket detection methods built on static structures have been
developed for decades. \citet{krone2016visual} classified the
mainstream approaches into four categories: grid-based,
Voronoi-based (e.g., \textsc{fpocket}), surface-based, and
probe-based methods. Tools such as $\text{EPOS}_{\text{BP}}$
\citep{eposbp_website} further enable high-throughput detection of
pocket ensembles on protein surfaces. Building on these algorithms,
a variety of pocket characterization and similarity metrics have been
proposed \citep{eguida2022estimating,weill2010alignment} to support
cross-protein comparison. The systematic evaluation of pocket
prediction methods by \citet{krivak2024comparative} shows that
geometry-based and machine-learning-based methods each have their
strengths, but both rely on a static structure as input.

\subsection{The Limitations of Static Pocket Detection}
Despite the substantial progress made in static pocket detection, its
fundamental limitation lies in \emph{neglecting the intrinsic
dynamics of protein conformations}. Real proteins undergo continuous
thermal motion, during which pocket shape, volume, accessibility, and
residue composition all vary over time. Detection based on a single
crystal structure (i)~\textbf{misses cryptic pockets} that appear
only in specific dynamic conformations \citep{meller2025cryptic};
(ii)~cannot characterize the \textbf{transient expansion,
contraction, merging, and splitting} of pockets
\citep{stank2016protein,lindow2013exploring}; and (iii)~struggles to
distinguish \textbf{functionally persistent pockets} from
\textbf{transient noise-like cavities} \citep{parulek2013visual}.
These limitations directly hinder the application of static pocket
analysis to allosteric regulation studies, drug repurposing, and the
targeting of ``undruggable'' proteins.

\subsection{The Importance of Dynamic Pocket Analysis}
With advances in hardware acceleration and enhanced sampling
techniques for molecular dynamics (MD) simulations, researchers can
now observe pocket evolution in conformational space at atomic
resolution \citep{stank2016protein,vant2024complex}. Dynamic pocket
analysis can reveal: the opening and closing mechanisms of cryptic
pockets \citep{meller2025cryptic}, transient accessibility along
ligand entry/exit channels \citep{lindow2013exploring}, and inter-site
allosteric communication networks
\citep{cuendet2021centrality,sheikamamuddy2021allosteric,
lasala2017pocket}. \citet{stank2016protein} systematically reviewed
multiple dimensions of pocket dynamics characterization --- shape,
volume, stability, and inter-pocket transitions --- and emphasized
their irreplaceable value in drug design. Recent reviews on cryptic
pockets further indicate that pocket discovery from a dynamic
perspective has become a key avenue for expanding the druggable
proteome \citep{meller2025cryptic}.

\subsection{Limitations of Alignment-Based Approaches for Dynamic
            Pocket Analysis}
In existing dynamic pocket analysis frameworks,
\emph{structural alignment} is an implicit prerequisite for most
methods: all frames of an MD trajectory must first be superimposed
onto a reference structure, and then pocket density, volume, or
shape can be aggregated and compared in a unified coordinate system
\citep{lindow2013exploring,parulek2013visual,zhan2020spatiotemporal}.
This paradigm suffers from several issues:
\begin{enumerate}
  \item \textbf{Reference selection is subjective} --- using
        different reference frames or different reference residue
        sets can substantially alter the resulting pocket density
        map \citep{krone2016visual}.
  \item \textbf{Alignment is unreliable in flexible regions} ---
        the motion of flexible loops and surface residues is often
        ``averaged out'' by global alignment, masking the true local
        pocket dynamics \citep{eguida2022estimating}.
  \item \textbf{Coupling between global and local motions} ---
        global translation and rotation of the protein cannot be
        cleanly separated from local pocket variation through
        alignment, and the pocket dynamics signal becomes
        contaminated by global pseudo-signals
        \citep{vant2024complex}.
  \item \textbf{Cross-protein comparison is difficult} --- when
        comparing pockets from different proteins or different
        mutants of the same protein, global alignment often fails
        because of differences in fold
        \citep{eguida2022estimating,weill2010alignment}.
\end{enumerate}

\subsection{The Advantages of Alignment-Free Approaches}
To circumvent the issues above, alignment-free approaches to pocket
analysis have attracted growing attention. \citet{weill2010alignment}
proposed an early ultra-high-throughput pocket comparison method
based on pharmacophoric fingerprints, demonstrating that
pocket-to-pocket similarity can be computed in milliseconds without
any structural alignment. In a 2022 review,
\citet{eguida2022estimating} further pointed out that alignment-free
methods offer significant advantages in large-scale pocket
comparison, cross-family comparison, and machine-learning-oriented
feature construction. The core advantages of alignment-free methods
are: (i)~inherent rotational and translational invariance, free from
reference-frame bias; (ii)~massive parallelizability, suitable for
trajectories of tens of thousands of frames; and (iii)~clear
biological semantics --- descriptions based on residue identity or
topological invariants naturally preserve cross-frame correspondence.
These advantages are particularly relevant to the cross-frame pocket
analysis on a single protein studied in this paper, because residue
numbering remains constant throughout the trajectory, making
alignment-free analysis essentially ``free of charge.''

\subsection{The Need for Visual Analytics on Long-Timescale MD Data}
As MD simulation timescales extend from nanoseconds into the
microsecond and even millisecond regime, a single trajectory may
contain tens of thousands to millions of frames covering hundreds
of thousands of atoms --- a high-dimensional spatiotemporal dataset
\citep{vant2024complex}. Numerical statistics alone (e.g., RMSD,
RMSF, pocket-volume time series) cannot reveal the complex events
embedded in pocket dynamics --- opening, closing, merging,
splitting, and allosteric coupling. \emph{Visual analytics}, by
combining automated algorithms with human perception and cognition,
provides a multi-layered, interactive, and exploratory path to
insight \citep{krone2016visual,parulek2013visual,vant2024complex}.
Both the state-of-the-art report by \citet{krone2016visual} and the
2024 review of MD trajectory visualization by
\citet{vant2024complex} emphasize that dedicated visual analytics
systems for dynamic pockets are a key bridge between raw MD data and
biological interpretation. The work presented in this paper is
motivated by this need: we propose a residue co-occurrence
network--based approach to dynamic pocket analysis and visualization.

% =====================================================================
\section{Related Work}
% =====================================================================

\subsection{Geometric and Topological Visual Analytics for
            Dynamic Pockets}
Early work on dynamic pocket visualization was primarily based on
geometric and topological representations.
\citet{lindow2013exploring} proposed a preprocessing pipeline based
on Voronoi diagrams and van der Waals spheres, coupled with an
interactive visualization interface that allows users to dynamically
select and observe cavity evolution along an MD trajectory.
\citet{parulek2013visual} proposed a pocket extraction method based
on implicit-function sampling and graph algorithms, used 3D graphs
and residue-based attributes to characterize pockets, and applied
the method to MD simulations of Proteinase~3.
\citet{krone2014visual}'s IEEE PacificVis 2014 work further proposed
a multiple-coordinated-view visual analytics system targeting dynamic
protein cavities and binding sites. The 2016 STAR by
\citet{krone2016visual} provided a comprehensive survey of the field
and proposed the four-class taxonomy of pocket detection methods, and
remains one of the most influential reviews in this area.

More refined visual analytics systems for dynamic pockets have
appeared in recent years: the spatiotemporal multiscale framework of
\citet{zhan2020spatiotemporal} presents pocket evolution at multiple
semantic levels; PNMAVis \citep{wang2022pnmavis} couples normal mode
analysis (NMA) with pocket dynamics to provide mode-based
interpretations of pocket variation; VAPPD \citep{vappd2022} proposes
a pocket shape representation that integrates topological features,
together with a high-dimensional similarity metric and a multi-view
analytical system; and the more recent PocketSCP \citep{pocketscp2025}
exploits the spatiotemporal topological properties of lining residue
atoms for visualizing and analyzing dynamic pockets. At the tool
level, D3Pockets \citep{d3pockets} provides a systematic web service
for pocket dynamics, while works such as
\textit{Visual Analysis of Large-Scale Protein--Ligand Interaction
Data} \citep{vis_largescale_pli} contribute complementary
perspectives from the angle of large-scale protein--ligand
interaction visualization.

\subsection{Residue Interaction Networks (RINs) and Their Use in MD}
Representing a protein structure as a residue interaction network
(RIN) has become an important paradigm for analyzing protein function
and dynamics \citep{vishveshwara2009modeling}. In this framework,
nodes represent residues and edges represent residue--residue
interactions in some defined sense, compressing high-dimensional
structural information into a graph that is amenable to network
analysis. Centrality measures computed on RINs (e.g., degree,
betweenness, closeness) have been widely used to identify
functionally important residues and allosteric hubs
\citep{cuendet2021centrality}.

As MD-derived data have become more abundant, researchers have
extended RINs to time-resolved settings.
\citet{rin_clustering_ensembles} extracted ensembles of RINs along
MD trajectories and clustered them, revealing distinct network
patterns that correspond to different conformational families. A
recent review by \citet{bernetti2024probing} systematically
summarizes methodological advances that combine MD with network
analysis for probing allosteric communication.
\citet{sheikamamuddy2021allosteric}, working on Falcipain~2, employed
dynamic RIN centrality analysis to identify allosteric pockets and
network hubs associated with drug-resistance mutations, providing a
representative example of how RINs can serve dynamic pocket analysis.

\subsection{Pocket-to-Pocket Communication and ``Pocket Crosstalk''
            Networks}
A few recent studies have started to construct networks at a higher
abstraction level than residues. \citet{lasala2017pocket} proposed
\emph{pocket crosstalk} analysis, in which the temporal exchange of
atoms between pockets along an MD trajectory is treated as the edge
of an allosteric communication network. By taking pockets as nodes
and inter-pocket atom flow as edges, this work complements
traditional residue-level networks at a higher hierarchical level
and reveals inter-pocket allosteric coupling patterns that had been
difficult to capture.

\subsection{Limitations of Existing Work and the Niche for This
            Paper}
The literature above can be seen as two relatively independent
threads: one focusing on \emph{geometry- and topology-driven dynamic
pocket visualization} (e.g., VAPPD, PocketSCP, the spatiotemporal
multiscale framework)
\citep{zhan2020spatiotemporal,vappd2022,pocketscp2025};
and the other focusing on \emph{network-driven residue analysis}
\citep{vishveshwara2009modeling,cuendet2021centrality,
rin_clustering_ensembles,bernetti2024probing}. However:

\begin{itemize}
  \item Existing residue interaction networks are largely built on
        physical distance or non-covalent interactions and lack a
        network representation at the level of \emph{functional
        co-membership};
  \item Existing dynamic pocket visualization approaches are mainly
        geometric or topological, and \emph{do not fully exploit the
        strengths of network analysis} for revealing pocket
        stability, merging/splitting events, and allosteric
        communication;
  \item Cross-frame pocket tracking has long depended on structural
        alignment, and the \emph{alignment-free paradigm has not yet
        been systematically adopted} in dynamic pocket analysis
        \citep{eguida2022estimating,weill2010alignment}.
\end{itemize}

To address these gaps, we propose the \textbf{Pocket-derived
Co-occurrence Residue Network (PCRN)}: treating the pockets detected
per-frame by \textsc{fpocket} as ``functional events,'' we use the
co-occurrence frequency of residue pairs within the same pocket as
edge weights and the per-residue pocket participation frequency as
node weights, thereby constructing a fully alignment-free,
pocket-oriented residue-level network. Combined with sliding time
windows and community detection, PCRN can characterize the stable
cores and dynamic edges of pockets, the merging and splitting events
along the trajectory, and the allosteric communication paths between
pockets, and presents these dynamic behaviors to domain experts
through a multi-view visual analytics system in an interpretable
manner.

% =====================================================================
\bibliography{references}
% =====================================================================

\end{document}
